\origpage{329}{N. Robertson, P.D. Seymour / Journal of Combinatorial Theory, Series B 90 (2004) 325--385}
\sectionhead{3. Frames and colour schemes}

In this section we cast 2.1 into a more convenient form. Let $\Sigma$ be a surface. A \textit{directed drawing in $\Sigma$} means a drawing in $\Sigma$ with a direction assigned to each edge. A \textit{frame $\Phi$ in $\Sigma$} consists of a directed drawing in $\Sigma$ (which we also denote by $\Phi$) with $U(\Phi)=bd(\Sigma)$, together with a designation of each edge of $\Phi$ as \textit{long} or \textit{short}, such that no two long edges have a common end. We call the edges of $\Phi$ the \textit{sides} of the frame. A \textit{painting $\Gamma$ in $\Sigma$} is said to \textit{fit a frame $\Phi$} if

\hspace*{12pt}$\bullet$ $V(\Phi)\subseteq V(\Gamma)$,

\hspace*{12pt}$\bullet$ each short side of $\Phi$ is an edge of $\Gamma$,

\hspace*{12pt}$\bullet$ for every border edge $e$ of $\Gamma$, if $S$ is the side of $\Phi$ with $e\subseteq S$ then the tail of $e$ precedes its head as $S$ is traversed from its tail in $\Phi$ to its head (that is, briefly, the direction of $e$ defined by $\gamma_\Gamma(e)$ agrees with the direction of $S$ in $\Phi$),

\hspace*{12pt}$\bullet$ $\Gamma$ is internally 3-connected,

\hspace*{12pt}$\bullet$ if $e\in E(\Gamma)$ and $|\tilde{e}|=3$ and $r$ is a region of $\Gamma$ in $\Sigma$ with $|\bar{r}\cap V(\Gamma)|\ge 3$, then $\bar{f}\cap bd(\Sigma)\ne \emptyset$ for every component $f$ of $e\cap \bar{r}$.

The fourth and fifth conditions have nothing to do with the frame, but this is a convenient place to introduce them. The fifth condition is a very strong requirement; in particular, it says that for any edge $e$ of $\Gamma$ with size 3, if no vertex of $\tilde{e}$ is in $bd(\Sigma)$, then for any two $u,v\in \tilde{e}$, there is a second edge $f$ with $u,v\in \tilde{f}$. This might seem to be very restrictive, but really it is not. For a given general painting, we can augment it to construct another satisfying this restriction (for every edge $e$ of size 3, just add three edges of size 2 joining the pairs of vertices of $\tilde{e}$, drawn close to $e$); and it turns out that one of these augmented paintings contains another in the required way, if and only if one of the unaugmented paintings contains another. So if we only want to prove well-quasi-ordering, it suffices to prove it for augmented paintings. And it turns out to be a convenient technical device for the proof, later.

If $e$ is a border edge of $\Gamma$ and $S$ is a side of $\Phi$ with $e\subseteq S$, we say that $e$ \textit{borders $S$}. If $\Gamma,\Gamma'$ are paintings in $\Sigma$ both fitting a frame $\Phi$, and $\sigma$ is an inflation of $\Gamma$ in $\Gamma'$, we say that $\sigma$ \textit{respects $\Phi$} if for every border edge $e$ of $\Gamma$, $\sigma(e)$ is a border edge and $e$ and $\sigma(e)$ border the same side of $\Phi$. If $\sigma$ respects $\Phi$, it follows that $\sigma(e)=e$ for every short side $e$ of $\Phi$, and $v\in V(\sigma(v))$ for every $v\in V(\Phi)$. If $\sigma$ respects $\Phi$, and satisfies the second condition in the definition of ``linear''\fix{, and maps every internal edge of $\Gamma$ to an internal edge of $\Gamma'$}, then $\sigma$ is linear\del{, as the reader should verify}. \fix{For the first condition, let $e \in E(\Gamma)$ and let $\Theta$ be a cuff. Each border edge of $\Gamma$ or of $\Gamma'$ is a connected subset of $bd(\Sigma) \setminus V(\Phi)$, since $V(\Phi)$ is a subset of $V(\Gamma)$ and of $V(\Gamma')$, and so it is a subset of some side of $\Phi$; and a side that meets $\Theta$ is a subset of $\Theta$, since it is connected. If $e$ borders $\Theta$, let $S$ be the side with $e \subseteq S$. Then $S \subseteq \Theta$, and $\sigma(e) \subseteq S$ since $\sigma$ respects $\Phi$; so $\sigma(e)$ borders $\Theta$. Conversely, suppose that $\sigma(e)$ borders $\Theta$. Then $\sigma(e)$ is a border edge, and so $e$ is a border edge, by the added hypothesis. Let $S$ be the side with $e \subseteq S$. Since $\sigma$ respects $\Phi$, $\sigma(e) \subseteq S$; so $S$ meets $\Theta$, and therefore $e \subseteq S \subseteq \Theta$, and $e$ borders $\Theta$. For the third condition, let $e$ border a cuff $\Theta$, let $S$ be the side with $e \subseteq S$, and orient $\Theta$ so that the tail of $e$ immediately precedes $e$. Since the direction of $e$ agrees with the direction of $S$, this orientation traverses $S$ from its tail to its head. Now $\sigma(e) \subseteq S$, and the direction of $\sigma(e)$ also agrees with the direction of $S$, since $\Gamma'$ fits $\Phi$; so the tail of $\sigma(e)$ immediately precedes $\sigma(e)$.}\fixnote{fix:3-linear}{As printed this is false. In a M\"obius band there are paintings $\Gamma, \Gamma'$ fitting a frame $\Phi$, and an inflation $\sigma$ of $\Gamma$ in $\Gamma'$ which respects $\Phi$ and satisfies the second condition, but maps an internal edge of $\Gamma$ to a border edge of $\Gamma'$, so that the first condition fails. The added hypothesis excludes this. The paper uses the remark only in the proof of 6.2, where the added hypothesis holds, as shown there. The argument above is the one used in the formalization.}

A \textit{colour scheme $\chi$} consists of a surface $\Sigma_\chi$, a frame $\Phi_\chi$ in $\Sigma_\chi$, two well-quasi-orders $\Omega_\chi(2)$ and $\Omega_\chi(3)$, and a well-quasi-order $\Omega_\chi(S)$ for each side $S$ of $\Phi_\chi$, such that $E(\Omega_\chi(S))=\{S\}$ for every short side $S$. (Thus, to specify a colour scheme, we need not define $\Omega_\chi(S)$ for short sides $S$.) If $\chi$ is a colour scheme, a \textit{$\chi$-coloured painting} is a pair $(\Gamma,\phi)$, where

\hspace*{12pt}$\bullet$ $\Gamma$ is a painting in $\Sigma_\chi$ fitting the frame $\Phi_\chi$,

\hspace*{12pt}$\bullet$ $\phi$ is a function with domain $E(\Gamma)$, such that if $e\in E(\Gamma)$ is internal then $\phi(e)\in \Omega_\chi(|\tilde{e}|)$, and if $e$ borders a side $S$ then $\phi(e)\in \Omega_\chi(S)$.
